\subsection{A two-sided recommendation benchmark}
\label{app:su-replication}
In the interview model of Su, Bayoumi and Joachims~\cite{rankings}, applicants may receive several invitations. Can we reproduce the congestion result before adapting their approach to rentals? We ran the authors' released solver and evaluation code with 30 applicants, 20 employers, crowding $0.5$ and inverse-rank attention. Each of ten market seeds (621--630) was simulated 10,000 times per policy. Naive, reciprocal and welfare ranking produced 18.41, 20.86 and 22.77 expected interviews, respectively. The release reports 18.7, 21.2 and 23.2. The authors give a default seed but do not list their full ten-seed set. We corrected an undefined variable used to save the final file after optimization.

We also evaluated expected interviews directly. If applicant $i$ applies to listing $j$ with probability $a_{ij}$, their expected examination probability given application is
\[
  \int_0^1 \prod_{h\succ_j i}(1-a_{hj}+a_{hj}t)\,dt.
\]
Indeed, if $N$ higher-priority applicants apply, $1/(1+N)=\int_0^1 t^N\,dt$. Independence across applicants gives the product. Multiplication by $a_{ij}$ and reciprocal-side relevance and summation over pairs gives expected interviews. With 100 employers and 150 applicants, the release-order welfare and reciprocal methods yielded 199.12 and 201.05 at zero crowding, 152.30 and 129.97 at half crowding and 117.37 and 91.33 at full crowding. The released optimizer iterates the inverse of the permutation of applicants in decreasing employer relevance. The equations specify the decreasing-relevance permutation itself. With this second order the corresponding welfare estimates were 203.52, 154.09 and 117.37. Both variants were evaluated using decreasing employer relevance, as in the authors' evaluator.

On the authors' 925-participant conference-networking matrix of imputed relevance, we removed self-pairs and reranked the first 100 reciprocal suggestions as specified in the paper. We obtained 603.0 and 764.1 expected successful replies for naive and reciprocal ranking, close to the published 604.0 and 763.9. Our equation-order welfare ranking reached 862.8, compared with the published 824.1. The policy ordering replicated. The difference in the optimized value remains unexplained. These outcomes are model expectations conditional on imputed relevance.

For the rental adaptation we used observable-score midranks within each renter's feasible listings and within each listing's eligible renters as proxies $f_{ij}$ and $g_{ij}$ for bilateral relevance. Each midrank is the average tied rank minus one half, divided by its comparison-set size. We used the paper's descending priority order. If $x_{ij}$ denotes display probability, application probability in the surrogate is $a_{ij}=f_{ij}x_{ij}$. We optimized its inverse-rank Jensen bound
\[
  \sum_{ij} \frac{g_{ij}a_{ij}}{1+\sum_{h\succ_j i}a_{hj}}
\]
over the row-wise convex hull of feasible top-five menus, initialized uniformly. Each of 50 Frank--Wolfe steps has size $0.2$. The linear oracle sorts the exposure gradient in each row. We displayed the five largest feasible final exposures (or all if fewer than five), resolving ties by independent market-keyed permutations. Deferred acceptance then ran on the resulting graph with unit capacities and the same latent preferences used for every housing comparison. The rank probabilities and the hard display quota are modeling choices. Expected replies and realized one-to-one allocations can favor different menus.
